\documentclass{article}
\usepackage[T1]{fontenc}
\usepackage{lmodern}
\usepackage{graphicx}
\usepackage{amsmath,amssymb}
\usepackage[a4paper,margin=2cm]{geometry}
\usepackage[absolute,overlay]{textpos}
\setlength{\TPHorizModule}{1mm}
\setlength{\TPVertModule}{1mm}
\usepackage[indonesian]{babel}
\usepackage{hyperref}
\setlength{\emergencystretch}{2em}
% unit: Halaman 44

\begin{document}

% === Halaman 44 (python/native) ===

• Pengisian kotak-S DES masih menjadi misteri. Nilai-nilai awal  S-box 

sudah diubah.

• Sebenarnya delapan putaran sudah cukup untuk membuat cipherteks 

sebagai fungsi acak dari setiap bit plainteks dan setiap bit kunci. 

• Namun dari penelitian dengan melakukan kriptanalisis, DES dengan 

jumlah putaran yang kurang dari 16 ternyata dapat dipecahkan 

dengan known-plaintext attack.

44

\end{document}
